\section{Conclusion}
\label{sec:conclusion}

This work asked how much of a reported tomato counting error belongs to the
measurement protocol and how much to the detector. Six protocol choices were
varied one at a time, and the training recipe was held fixed wherever the
question allowed it. The choices were the annotation convention that defines a
maturity stage, the occluder, the confidence threshold, the illumination of the
test frames, the model size, and the runtime used to execute the model.

The measured effects differ in magnitude, and model size does not produce the
largest change. Across the YOLO11 ladder trained under one protocol, clean
recall spans 0.058 and mAP50--95 spans 0.068. Both stop improving at YOLO11m,
while package size, latency, and load time continue to grow through YOLO11x.
Over-exposing the same validation frames by three stops reduces the strongest
model's recall by 0.205. A cool white-balance cast reduces recall by 0.05 to
0.08. Hiding 40\% of each fruit reduces recall by more than 0.3. This result
also depends on the occluder's appearance: recall is 0.571 under a same-image
background patch and 0.461 under a black patch at the same hidden fraction.
Changing the annotation convention places 98\% of one dataset's fully ripe
boxes within the other dataset's half-ripened band. The class-level transfer
error follows the same direction. Changing the confidence threshold alters
cross-domain counting error by more than a factor of two while leaving mAP
unchanged. Selecting the threshold separately for each model also reorders the
counting results in the size ladder. The FP16 Core ML export changes recall by
at most 0.002. For the same package, latency varies with the compute units
(2.0$\times$ to 4.1$\times$), the device (1.21$\times$ to 3.85$\times$ against
a laptop), the operating-system version (67\% for the smallest model across a
single update), and the run length (8.8\% drift over 4{,}500 frames for the
largest).

A protocol based on one summary number would have changed three conclusions.
An ablation evaluated on clean held-out frames would have recommended removing
the colour jitter. Removing it has no measured cost on clean data but reduces
recall by up to 0.102 after perturbation. A single training run at the most
perturbed level reported the default recipe at the most favourable of its three
seeds. The tuned counting column of the size ladder implies that counting
improves with capacity, while the fixed-threshold column does not support that
conclusion.

These results establish a reporting requirement. A reproducible counting error
must identify the annotation source and the colour boundary used by that
source. An occlusion result must identify the occluder and matching rule. A
counting result must report the confidence threshold and the split used to
select it, while a lighting evaluation must state its illumination range. A
latency result must identify the device, compute units, operating system, and
run duration. Reporting these details separately allows the detection, staging,
counting, and deployment results to be reproduced. A single aggregate number
does not preserve this information.

For the application that motivated the work, the practical conclusions are
specific. YOLO11n and YOLO11s meet a 30\,FPS camera budget on the smartphone's
Neural Engine with three to four times the required margin at their p90
latency. Models above YOLO11m do not improve accuracy, so additional capacity
does not address the measured limitations of this system. The default HSV
jitter should be retained because removing it costs an order of magnitude more
than widening it returns. Fully-ripe counting error should be reported alongside
mAP because the two metrics do not follow the same pattern. The confidence
threshold should be selected on frames that are not used to report the result.

Future work should test the photometric results on a collection recaptured under
measured illumination. The current sweep is a simulation applied to existing
frames, so the additional collection would show whether the ordering holds with
real sensor behaviour. Instance masks would improve the box-overlap estimate of
natural occlusion, although they would still not establish depth order. Seed
repeats across the whole ladder, beyond the colour recipes, would allow the
differences among YOLO11m, l, and x to be ranked instead of reported as
indistinguishable. The deployment study also lacks a comparison between a
training-time intervention and exposure control at capture time. Such a
comparison could determine whether capture control recovers the over-exposure
loss at a lower cost than changing the model.
